\documentclass[a4paper,11pt]{article}

\usepackage[T1]{fontenc}
\usepackage[utf8]{inputenc}
\usepackage[french]{babel}
\usepackage{geometry}
\usepackage{xcolor}
\usepackage{titlesec}
\usepackage{listings}
\usepackage{fancyhdr}

\geometry{margin=2.2cm}
\emergencystretch=3em
\sloppy

\definecolor{alanyaBlue}{RGB}{40,60,120}
\definecolor{codebg}{RGB}{246,246,248}
\definecolor{codeframe}{RGB}{210,210,216}
\definecolor{sqlkw}{RGB}{40,60,120}
\definecolor{sqlcomment}{RGB}{120,120,120}

\titleformat{\section}{\large\bfseries\color{alanyaBlue}}{Optimisation \thesection}{0.5em}{\ ---\ }
\titlespacing*{\section}{0pt}{1.3em}{0.5em}

\pagestyle{fancy}
\fancyhf{}
\renewcommand{\headrulewidth}{0pt}
\fancyfoot[C]{\small\thepage}

\lstdefinelanguage{SQLplus}{
  language=SQL,
  morekeywords={AUTO_INCREMENT, ENGINE, INNODB, UNSIGNED, IGNORE, GENERATED, ALWAYS, STORED, CONSTRAINT},
}
\lstset{
  language=SQLplus,
  basicstyle=\ttfamily\footnotesize,
  keywordstyle=\color{sqlkw}\bfseries,
  commentstyle=\color{sqlcomment}\itshape,
  backgroundcolor=\color{codebg},
  frame=single,
  rulecolor=\color{codeframe},
  framesep=6pt,
  xleftmargin=4pt,
  xrightmargin=4pt,
  breaklines=true,
  breakatwhitespace=false,
  showstringspaces=false,
  columns=fullflexible,
  tabsize=2,
}

\title{\color{alanyaBlue}\textbf{ALANYA Optimisations de la base de données}\\}
\date{24 Août 2026}
 
\begin{document}
\maketitle
\thispagestyle{fancy} 
\section{Index sur les chemins de lecture/écriture les plus chauds}

\begin{lstlisting}
ALTER TABLE message ADD INDEX idx_message_conv_status
  (conversationID, status, senderID);

ALTER TABLE message ADD INDEX idx_message_conv_msgid
  (conversationID, msgID);

ALTER TABLE job_queue ADD INDEX idx_job_ready
  (failed_at, locked_at, run_after, id);

ALTER TABLE users ADD INDEX idx_users_fcm_token (fcm_token(191));

ALTER TABLE users DROP INDEX idx_users_phone;
\end{lstlisting}

\noindent\textit{Commentaire.} Les deux premiers index évitent un balayage complet de l'historique d'une conversation à chaque accusé de lecture et à chaque synchronisation delta. Le troisième supprime un balayage complet de la file de tâches à chaque cycle de traitement (jusqu'à 20 fois par seconde). Le dernier \texttt{DROP} retire un index strictement redondant avec la contrainte \texttt{UNIQUE} existante.
 

\section{Élargissement de \texttt{statut.ID} en BIGINT}

\begin{lstlisting}
ALTER TABLE statut_views            DROP FOREIGN KEY fk_sv_statut;
ALTER TABLE statut_i18n             DROP FOREIGN KEY fk_statut_i18n;
ALTER TABLE welcome_status_delivery DROP FOREIGN KEY fk_wsd_statut;

ALTER TABLE statut MODIFY ID BIGINT UNSIGNED NOT NULL AUTO_INCREMENT;
ALTER TABLE statut_views MODIFY statutID BIGINT UNSIGNED NOT NULL;
ALTER TABLE statut_i18n MODIFY statut_id BIGINT UNSIGNED NOT NULL;
ALTER TABLE welcome_status_delivery MODIFY statut_id BIGINT UNSIGNED NOT NULL;
ALTER TABLE broadcast MODIFY statut_id BIGINT UNSIGNED NULL;

ALTER TABLE statut_views ADD CONSTRAINT fk_sv_statut
  FOREIGN KEY (statutID) REFERENCES statut(ID) ON DELETE CASCADE;
ALTER TABLE statut_i18n ADD CONSTRAINT fk_statut_i18n
  FOREIGN KEY (statut_id) REFERENCES statut(ID) ON DELETE CASCADE;
ALTER TABLE welcome_status_delivery ADD CONSTRAINT fk_wsd_statut
  FOREIGN KEY (statut_id) REFERENCES statut(ID) ON DELETE CASCADE;
\end{lstlisting}

\noindent\textit{Commentaire.} \texttt{statut.ID} était en \texttt{INT} signé (plafond 2,1~milliards), saturable en une quinzaine de mois au rythme visé. Les trois clés étrangères référençantes doivent être déposées avant l'élargissement de la colonne référencée, puis recréées à l'identique.

\section{Élargissement des colonnes sur les tables qui grossiront rapidement}

\begin{lstlisting}
ALTER TABLE broadcast_delivery
  MODIFY conversID BIGINT UNSIGNED NULL,
  MODIFY msgID     BIGINT UNSIGNED NULL;

ALTER TABLE welcome_delivery
  MODIFY conversID BIGINT UNSIGNED NOT NULL;

ALTER TABLE trip_watcher
  MODIFY msgID     BIGINT UNSIGNED NULL,
  MODIFY conversID BIGINT UNSIGNED NULL;

ALTER TABLE conv_participants
  MODIFY pendingJoinMsgID BIGINT UNSIGNED NULL,
  MODIFY unreadCount      INT UNSIGNED NOT NULL DEFAULT 0;
\end{lstlisting}

\noindent\textit{Commentaire.} Ces colonnes référençaient des identifiants \texttt{BIGINT} (\texttt{conversation.conversID}, \texttt{message.msgID}) tout en étant elles-mêmes en \texttt{INT} ou en type désaccordé, avec un risque de troncature silencieuse au-delà du plafond. \texttt{unreadCount} était en \texttt{SMALLINT} signé (plafond 32\,767), débordable sur un groupe très actif.
%
%\section{Purges automatiques de rétention}
%
%\begin{lstlisting}
%-- Index de support (une purge sans index scanne sa table entiere)
%ALTER TABLE appareils          ADD INDEX idx_appareils_revoked (revoked_at);
%ALTER TABLE user_push_devices  ADD INDEX idx_push_heartbeat (lastHeartbeatAt);
%ALTER TABLE users               ADD INDEX idx_users_reset_otp_exp (reset_otp_expires_at);
%ALTER TABLE users               ADD INDEX idx_users_email_otp_exp (email_change_otp_expires_at);
%ALTER TABLE broadcast_delivery  ADD INDEX idx_delivery_purge (delivered_at);
%ALTER TABLE user_export_jobs    ADD INDEX idx_export_jobs_expires (expiresAt);
%
%-- Requetes executees quotidiennement par le service de retention
%DELETE FROM statut WHERE expiredAt < DATE_SUB(NOW(), INTERVAL 7 DAY) LIMIT 5000;
%
%DELETE FROM callHistory WHERE created_at < DATE_SUB(NOW(), INTERVAL 12 MONTH) LIMIT 5000;
%
%DELETE FROM userAccess WHERE dateLogin < DATE_SUB(NOW(), INTERVAL 90 DAY) LIMIT 5000;
%
%DELETE FROM job_queue
%  WHERE failed_at IS NOT NULL AND failed_at < DATE_SUB(NOW(), INTERVAL 30 DAY)
%  LIMIT 5000;
%
%DELETE FROM appareils
%  WHERE revoked_at IS NOT NULL AND revoked_at < DATE_SUB(NOW(), INTERVAL 90 DAY)
%  LIMIT 1000;
%
%DELETE FROM user_push_devices
%  WHERE lastHeartbeatAt IS NOT NULL AND lastHeartbeatAt < DATE_SUB(NOW(), INTERVAL 180 DAY)
%  LIMIT 1000;
%
%UPDATE users SET reset_otp = NULL, reset_otp_expires_at = NULL
%  WHERE reset_otp_expires_at IS NOT NULL AND reset_otp_expires_at < NOW();
%
%UPDATE users SET email_change_otp = NULL, email_change_otp_expires_at = NULL,
%                  pending_email = NULL
%  WHERE email_change_otp_expires_at IS NOT NULL AND email_change_otp_expires_at < NOW();
%\end{lstlisting}
%
%\noindent\textit{Commentaire.} Aucune table à croissance non bornée n'était purgée avant ce travail. Huit tables sont désormais couvertes par un service exécuté quotidiennement sous bail, chaque requête plafonnée par \texttt{LIMIT} pour ne jamais tenir un verrou long sur un arriéré important.

\section{Colonne générée indexée pour les réactions}

\begin{lstlisting}
ALTER TABLE message
  ADD COLUMN has_reactions TINYINT GENERATED ALWAYS AS
    (CASE WHEN reactions IS NOT NULL AND JSON_LENGTH(reactions) > 0
          THEN 1 ELSE 0 END) STORED,
  ADD INDEX idx_message_conv_hasreactions (conversationID, has_reactions);
\end{lstlisting}

\noindent\textit{Commentaire.} Remplace un filtre \texttt{JSON\_LENGTH(reactions) > 0} non indexable, qui provoquait un balayage complet de l'historique d'une conversation à chaque ouverture d'écran. La colonne \texttt{STORED} se recalcule automatiquement à chaque écriture, sans changement côté code d'écriture. Bien que l'idéal serait une table dédiée pour les réactions.

\section{Extraction de \texttt{message.mediaThumb}}

\begin{lstlisting}
CREATE TABLE message_thumb (
  msgID BIGINT NOT NULL PRIMARY KEY,
  thumb MEDIUMBLOB NOT NULL,
  CONSTRAINT fk_message_thumb FOREIGN KEY (msgID)
    REFERENCES message(msgID) ON DELETE CASCADE
) ENGINE=InnoDB DEFAULT CHARSET=utf8mb4;

INSERT INTO message_thumb (msgID, thumb)
SELECT msgID, FROM_BASE64(mediaThumb) FROM message
WHERE mediaThumb IS NOT NULL AND mediaThumb <> '';

-- En attente (periode d'observation) :
-- ALTER TABLE message DROP COLUMN mediaThumb;
\end{lstlisting}

\noindent\textit{Commentaire.} La vignette base64 (mesurée à 36~Ko en moyenne) était chargée par tous les \texttt{SELECT m.*} sur la table la plus utilisée du schéma, utile ou non au chemin concerné. Stockage en binaire (\texttt{MEDIUMBLOB}) plutôt qu'en base64 : environ 25~\% d'économie. Le contrat JSON envoyé au client ne change pas (recodage à la lecture via \texttt{TO\_BASE64}). Le \texttt{DROP COLUMN} final reste volontairement en attente comme filet de rollback.

\section{Extraction de \texttt{is\_online}/\texttt{last\_seen} vers \texttt{user\_presence}}

\begin{lstlisting}
CREATE TABLE user_presence (
  alanyaID  INT      NOT NULL PRIMARY KEY,
  is_online TINYINT  NOT NULL DEFAULT 0,
  last_seen DATETIME NOT NULL DEFAULT CURRENT_TIMESTAMP,
  CONSTRAINT fk_user_presence FOREIGN KEY (alanyaID)
    REFERENCES users(alanyaID) ON DELETE CASCADE
) ENGINE=InnoDB DEFAULT CHARSET=utf8mb4;

INSERT INTO user_presence (alanyaID, is_online, last_seen)
SELECT alanyaID, is_online, last_seen FROM users;

-- En attente (periode d'observation) :
-- ALTER TABLE users DROP COLUMN is_online, DROP COLUMN last_seen;
\end{lstlisting}

\noindent\textit{Commentaire.} Ces deux colonnes, écrites à chaque connexion/déconnexion/bascule premier plan, entraient en contention directe avec le verrou exclusif posé par \texttt{materializeForUser} sur la ligne \texttt{users}. Les extraire supprime cette contention sans changer la stratégie de verrouillage. Bascule complète des lectures/écritures effectuée côté application (4~sites d'écriture, une quinzaine de requêtes de lecture) ; le \texttt{DROP COLUMN} final reste en attente.

\section{Extraction des colonnes de chiffrement vers \texttt{message\_crypto}} ( À utiliser plus tard pour la crypto)

\begin{lstlisting}
CREATE TABLE message_crypto (
  msgID                BIGINT        NOT NULL PRIMARY KEY,
  ciphertext           MEDIUMBLOB    NULL,
  dr_nonce             VARBINARY(16) NULL,
  dr_header            TEXT          NULL,
  archive_blob         MEDIUMBLOB    NULL,
  signal_message_type  TINYINT       NULL,
  CONSTRAINT fk_message_crypto FOREIGN KEY (msgID)
    REFERENCES message(msgID) ON DELETE CASCADE
) ENGINE=InnoDB DEFAULT CHARSET=utf8mb4;

INSERT INTO message_crypto
  (msgID, ciphertext, dr_nonce, dr_header, archive_blob, signal_message_type)
SELECT msgID, ciphertext, dr_nonce, dr_header, archive_blob, signal_message_type
FROM message
WHERE ciphertext IS NOT NULL OR dr_nonce IS NOT NULL OR dr_header IS NOT NULL
   OR archive_blob IS NOT NULL OR signal_message_type IS NOT NULL;

ALTER TABLE message
  DROP COLUMN ciphertext,
  DROP COLUMN dr_nonce,
  DROP COLUMN dr_header,
  DROP COLUMN archive_blob,
  DROP COLUMN signal_message_type;
\end{lstlisting}

\noindent\textit{Commentaire.} Cinq colonnes issues d'une fonctionnalité de chiffrement bout-en-bout en pause. La nouvelle table sera utilisée lors de la crypto

\section{Reclustering de \texttt{conv\_participants}}

\begin{lstlisting}
ALTER TABLE conv_participants
  DROP PRIMARY KEY,
  DROP INDEX uq_conv_user,
  ADD PRIMARY KEY (conversID, alanyaID),
  ADD UNIQUE KEY uq_cp_id (id);
\end{lstlisting}

\noindent\textit{Commentaire.} InnoDB organise physiquement une table selon sa clé primaire. Avec un identifiant synthétique, les participants d'une même conversation étaient dispersés sur des pages disque différentes. La nouvelle clé primaire \texttt{(conversID, alanyaID)} les colocalise, au bénéfice du chemin le plus sollicité de cette table : le fan-out des messages à tous les participants d'une conversation. L'identifiant \texttt{id} est conservé en clé secondaire (toujours \texttt{AUTO\_INCREMENT}).

\end{document}
